% GENERATED FILE. Edit canonical lessons, then run npm run generate:books.
% canonical-source-hash: fnv1a64:2b5389ecf5b786e8
% canonical-lessons: LA-C497-sollicitudo, LA-C497-ango, LA-C497-aegritudo, LA-C497-molestia, LA-C497-crucio

\chapter{Worry and distress}
\label{ch:la-worry-and-distress}
\hlchaptermodality{497}

\begin{chapteropening}
\textbf{By the end of this chapter:} \emph{I can talk about worry and distress.}

Complete the last lesson of chapter 497: I can talk about worry and distress.
\end{chapteropening}

\section[sollicitūdō, sollicitūdinis]{sollicitūdō, sollicitūdinis — anxiety, worry}

\label{lesson:LA-C497-sollicitudo}



\begin{quote}
Before the new one: say the Latin for dreadful, then the Latin for frightful.
\end{quote}

\subsection*{You'll want to know: sollicitūdō, sollicitūdinis}
\textbf{sollicitūdō, sollicitūdinis} — \textquotedblleft{}anxiety, worry\textquotedblright{}.

\begin{quote}
Māter magnā sollicitūdine fīlium exspectat.
\end{quote}

\begin{quote}
\emph{The mother waits for her son with great anxiety.}
\end{quote}

Feminine, third declension, from \emph{sollicitus} — anxious. English \emph{solicitude}.

\begin{cousinweb}
\textbf{sollicitūdō} is built on \emph{sollicitus} \textquotedblleft{}anxious, troubled\textquotedblright{} with \emph{-tūdō} (Lewis \& Short; EtymDB). English \emph{solicitude} comes from it.
\end{cousinweb}

\begin{grammarlens}[title={What you've built}]
One more word for telling, explaining and describing at B1.
\end{grammarlens}

\subsection*{Your turn}
\begin{itemize}
\raggedright
  \item \emph{Say it:} \emph{sollicitūdō, sollicitūdinis}
  \item \emph{Say it:} \emph{sollicitūdō, sollicitūdinis}, once more
  \item \emph{Say it:} a sentence of your own with \emph{sollicitūdō, sollicitūdinis}
  \item \emph{Recall:} say the Latin for dreadful, then the Latin for frightful, then say \emph{sollicitūdō, sollicitūdinis} again
\end{itemize}

\begin{tcolorbox}[breakable,colback=teal!4,colframe=teal!35!black,title={Before you move on}]
What does sollicitūdō, sollicitūdinis mean? (\textquotedblleft{}Anxiety, worry\textquotedblright{}.) Say it once more.
\end{tcolorbox}
\section[angō, angere]{angō, angere — to distress, torment; to choke}

\label{lesson:LA-C497-ango}



\begin{quote}
Before the new one: say the Latin for frightful, then the Latin for anxiety.
\end{quote}

\subsection*{You'll want to know: angō, angere}
\textbf{angō, angere} — \textquotedblleft{}to distress, torment; to choke\textquotedblright{}.

\begin{quote}
Cūra dē fīliō patrem angit.
\end{quote}

\begin{quote}
\emph{Worry about his son torments the father.}
\end{quote}

\emph{angō, angere, anxī}. First to squeeze the throat, then to torment the mind. \emph{Anxius}, already known, comes from it.

\begin{cousinweb}
\textbf{angō} is inherited from Proto-Italic \emph{angō}, from a Proto-Indo-European root meaning \textquotedblleft{}narrow, tight\textquotedblright{} (EtymDB). English \emph{anxious} and \emph{anguish} come from its family.
\end{cousinweb}

\begin{grammarlens}[title={What you've built}]
One more word for telling, explaining and describing at B1.
\end{grammarlens}

\subsection*{Your turn}
\begin{itemize}
\raggedright
  \item \emph{Say it:} \emph{angō, angere}
  \item \emph{Say it:} \emph{angō, angere}, once more
  \item \emph{Say it:} a sentence of your own with \emph{angō, angere}
  \item \emph{Recall:} say the Latin for frightful, then the Latin for anxiety, then say \emph{angō, angere} again
\end{itemize}

\begin{tcolorbox}[breakable,colback=teal!4,colframe=teal!35!black,title={Before you move on}]
What does angō, angere mean? (\textquotedblleft{}To distress, torment; to choke\textquotedblright{}.) Say it once more.
\end{tcolorbox}
\section[aegritūdō, aegritūdinis]{aegritūdō, aegritūdinis — distress of mind, grief; sickness}

\label{lesson:LA-C497-aegritudo}



\begin{quote}
Before the new one: say the Latin for anxiety, then the Latin for to distress.
\end{quote}

\subsection*{You'll want to know: aegritūdō, aegritūdinis}
\textbf{aegritūdō, aegritūdinis} — \textquotedblleft{}distress of mind, grief; sickness\textquotedblright{}.

\begin{quote}
Post mortem amīcī aegritūdō eum diū tenet.
\end{quote}

\begin{quote}
\emph{After his friend's death, grief holds him for a long time.}
\end{quote}

Feminine, third declension, from \emph{aeger} — sick. In philosophy it often means sorrow, a sickness of the mind.

\begin{cousinweb}
\textbf{aegritūdō} is built on \emph{aeger} \textquotedblleft{}sick\textquotedblright{} with \emph{-tūdō} (Lewis \& Short), and it names sickness of body or of mind.
\end{cousinweb}

\begin{grammarlens}[title={What you've built}]
One more word for telling, explaining and describing at B1.
\end{grammarlens}

\subsection*{Your turn}
\begin{itemize}
\raggedright
  \item \emph{Say it:} \emph{aegritūdō, aegritūdinis}
  \item \emph{Say it:} \emph{aegritūdō, aegritūdinis}, once more
  \item \emph{Say it:} a sentence of your own with \emph{aegritūdō, aegritūdinis}
  \item \emph{Recall:} say the Latin for anxiety, then the Latin for to distress, then say \emph{aegritūdō, aegritūdinis} again
\end{itemize}

\begin{tcolorbox}[breakable,colback=teal!4,colframe=teal!35!black,title={Before you move on}]
What does aegritūdō, aegritūdinis mean? (\textquotedblleft{}Distress of mind, grief; sickness\textquotedblright{}.) Say it once more.
\end{tcolorbox}
\section[molestia, molestiae]{molestia, molestiae — annoyance, trouble}

\label{lesson:LA-C497-molestia}



\begin{quote}
Before the new one: say the Latin for to distress, then the Latin for distress of mind.
\end{quote}

\subsection*{You'll want to know: molestia, molestiae}
\textbf{molestia, molestiae} — \textquotedblleft{}annoyance, trouble\textquotedblright{}.

\begin{quote}
Muscae aegrō molestiam afferunt.
\end{quote}

\begin{quote}
\emph{The flies cause the sick man annoyance.}
\end{quote}

Feminine, first declension, from \emph{molestus} — annoying. English \emph{molest} comes from the related verb \emph{molestō}.

\begin{cousinweb}
\textbf{molestia} is built on \emph{molestus} \textquotedblleft{}troublesome\textquotedblright{}, which Lewis \& Short connects with \emph{mōlēs} \textquotedblleft{}a heavy mass\textquotedblright{}: a burden. English \emph{molest} comes from the related verb \emph{molestō}.
\end{cousinweb}

\begin{grammarlens}[title={What you've built}]
One more word for telling, explaining and describing at B1.
\end{grammarlens}

\subsection*{Your turn}
\begin{itemize}
\raggedright
  \item \emph{Say it:} \emph{molestia, molestiae}
  \item \emph{Say it:} \emph{molestia, molestiae}, once more
  \item \emph{Say it:} a sentence of your own with \emph{molestia, molestiae}
  \item \emph{Recall:} say the Latin for to distress, then the Latin for distress of mind, then say \emph{molestia, molestiae} again
\end{itemize}

\begin{tcolorbox}[breakable,colback=teal!4,colframe=teal!35!black,title={Before you move on}]
What does molestia, molestiae mean? (\textquotedblleft{}Annoyance, trouble\textquotedblright{}.) Say it once more.
\end{tcolorbox}
\section[cruciō, cruciāre]{cruciō, cruciāre — to torment, torture}

\label{lesson:LA-C497-crucio}



\begin{quote}
Before the new one: say the Latin for distress of mind, then the Latin for annoyance.
\end{quote}

\subsection*{You'll want to know: cruciō, cruciāre}
\textbf{cruciō, cruciāre} — \textquotedblleft{}to torment, torture\textquotedblright{}.

\begin{quote}
Dolor dentium senem cruciat.
\end{quote}

\begin{quote}
\emph{Toothache torments the old man.}
\end{quote}

\emph{cruciō, cruciāre, cruciāvī, cruciātum}, from \emph{crux} — a cross. In comedy also of worry: \emph{mē cruciō} — I fret. English \emph{excruciating}.

\begin{cousinweb}
\textbf{cruciō} is built on \emph{crux} \textquotedblleft{}a cross\textquotedblright{} (Lewis \& Short; EtymDB): to torment as on a cross. English \emph{excruciating} comes from the compound \emph{excruciō}.
\end{cousinweb}

\begin{grammarlens}[title={What you've built}]
One more word for telling, explaining and describing at B1.
\end{grammarlens}

\subsection*{Your turn}
\begin{itemize}
\raggedright
  \item \emph{Say it:} \emph{cruciō, cruciāre}
  \item \emph{Say it:} \emph{cruciō, cruciāre}, once more
  \item \emph{Say it:} a sentence of your own with \emph{cruciō, cruciāre}
  \item \emph{Recall:} say the Latin for distress of mind, then the Latin for annoyance, then say \emph{cruciō, cruciāre} again
\end{itemize}

\begin{tcolorbox}[breakable,colback=teal!4,colframe=teal!35!black,title={Before you move on}]
What does cruciō, cruciāre mean? (\textquotedblleft{}To torment, torture\textquotedblright{}.) Say it once more.
\end{tcolorbox}
